\documentclass[11pt]{article}
\usepackage[T1]{fontenc}
\usepackage{lmodern}
\usepackage[margin=1in]{geometry}
\usepackage{amsmath,amssymb,amsthm,booktabs}
\usepackage[hidelinks]{hyperref}
\newtheorem{theorem}{Theorem}
\newtheorem{lemma}[theorem]{Lemma}
\DeclareMathOperator{\sgn}{sgn}
\newcommand{\E}{\mathbb E}
\newcommand{\Prb}{\mathbb P}
\title{Two asymptotic formulas for square Hardinian arrays}
\author{}
\date{October 13, 2026}
\begin{document}
\maketitle

\begin{abstract}
We prove that the numbers of square Hardinian arrays with parameters
$r=2$ and $r=3$ satisfy
$H_2(n,n)\sim16^n/(324\pi n)$ and
$H_3(n,n)\sim64^n/(78732\pi n^3)$, respectively.
These are two of the asymptotic formulas recorded as conjectures by Dougherty-Bliss
and Kauers. Starting from their complementary-minor counting identity,
we obtain exact sums of squares of binomial sums. A uniform local
expansion, together with a geometric tail bound, reduces the leading
constants to two small moment determinants. The argument treats fixed
$r=2,3$ as $n\to\infty$ and does not use a conjectured recurrence.
\end{abstract}

\section{Statement and source identity}

Following Dougherty-Bliss and Kauers~\cite{dk}, let $H_r(n,n)$ be the
number of arrays $(a_{i,j})_{1\le i,j\le n}$ of nonnegative integers such
that
\begin{enumerate}
\item $a_{1,1}=0$ and $a_{n,n}=n-r-1$;
\item every step one position to the right, down, or diagonally down and
right increases the entry by either $0$ or $1$;
\item $|a_{i,j}-(\max(i,j)-1)|\le r$ for all $i,j$.
\end{enumerate}
The step condition is imposed whenever both positions belong to the
array. We consider $r\in\{2,3\}$ fixed.

\begin{theorem}\label{thm:main}
As $n\to\infty$ through the positive integers,
\[
 H_2(n,n)\sim\frac{16^n}{324\pi n},\qquad
 H_3(n,n)\sim\frac{64^n}{78732\pi n^3}.
\]
\end{theorem}

The two formulas occur in the conjectural asymptotic table on page~12
of the published version of~\cite{dk}. OEIS A253217 credits the $r=2$
asymptotic conjecture to Vaclav Kotesovec on March 2, 2023~\cite{oeis2}.
Our starting point is the
counting identity displayed in the proof of its Theorem~6, on page~9;
we use the identity itself, rather than its consequence of
D-finiteness. Put $N=n-2$ and let
\[
 M=\left[\binom{i+j}{i}\right]_{0\le i,j\le N}.
\]
For $n\ge r+1$, that identity reads
\begin{equation}\label{eq:source}
 H_r(n,n)=\sum_{\substack{I,J\subseteq\{0,\ldots,N\}\\|I|=|J|=r}}
              \det M_{I^c,J^c}.
\end{equation}
Rows and columns of every submatrix are placed in increasing order;
complements are taken in $\{0,\ldots,N\}$ and the empty determinant is
$1$. Thus the matrix has $n-1$ rows and columns, and exactly $r$ of each
are deleted. In particular, each subset occurs once, with no factorial
factor. At $n=r+1$, both the array definition and~\eqref{eq:source}
give $1$.

We derive the theorem from~\eqref{eq:source}. We make no claim here
about the asymptotic formulas for general $r$, or about completion of
an exhaustive priority search.

\section{A sum of squares}

Let $P_{u,i}=\binom ui$ for $0\le u,i\le N$, with a binomial
coefficient equal to zero outside its usual range. The Vandermonde
identity gives $M=PP^T$. The matrix $P$ is unit lower triangular, and
its inverse has entries
\[
 (P^{-1})_{u,i}=(-1)^{u+i}\binom ui.
\]
Writing $E=\operatorname{diag}((-1)^0,\ldots,(-1)^N)$ and $B=P^TP$,
we therefore have $\det M=1$ and $M^{-1}=EBE$.
Jacobi's complementary-minor identity now gives
\[
 \det M_{I^c,J^c}
 =(-1)^{\sum I+\sum J}\det(M^{-1})_{J,I}
 =\det B_{J,I}.
\]
The last equality uses the identical sign supplied by the two diagonal
matrices $E$. Applying Cauchy--Binet to $B=P^TP$ and summing over $I,J$
in~\eqref{eq:source} yields
\begin{equation}\label{eq:squares}
 H_r(n,n)=\sum_{0\le u_1<\cdots<u_r\le N}W_r(u_1,\ldots,u_r)^2,
\end{equation}
where
\begin{equation}\label{eq:W}
 W_r(u_1,\ldots,u_r)=
 \sum_{0\le i_1<\cdots<i_r\le N}
 \det\left[\binom{u_a}{i_b}\right]_{a,b=1}^r.
\end{equation}
Indeed, for each row set $U$, the sums over $I$ and $J$ are the same
sum of minors, so their product is its square. All sums in this step
are finite.

For independent random variables $X_a\sim\operatorname{Bin}(u_a,1/2)$,
define
\[
 w_r(u_1,\ldots,u_r)=2^{-\sum_a u_a}W_r(u_1,\ldots,u_r).
\]
Expanding the determinant in~\eqref{eq:W} and sorting the distinct
values of the $X_a$ shows that
\begin{equation}\label{eq:prob}
 w_r(u_1,\ldots,u_r)
 =\E\!\left[\sgn\prod_{a<b}(X_b-X_a)\right].
\end{equation}
For distinct values, the determinant sign is exactly the sign of the
sorting permutation. A tie contributes zero on the right and is
absent from the strictly increasing column sets on the left. Here
$\sgn(0)=0$. In particular, we have the global bound
\begin{equation}\label{eq:globalbound}
 |w_r(u_1,\ldots,u_r)|\le1.
\end{equation}

\section{Exact formulas for two and three rows}

For integers $0\le u<v$, set
\begin{equation}\label{eq:D}
 D(u,v)=\sum_{j=u}^{v-1}\binom{u+v}{j},\qquad
 Q(u,v)=2^{-u-v}D(u,v).
\end{equation}
If $X\sim\operatorname{Bin}(u,1/2)$ and
$Y\sim\operatorname{Bin}(v,1/2)$ are independent, then $Y-X$ has
the distribution of $Z-u$ with $Z\sim\operatorname{Bin}(u+v,1/2)$.
By symmetry, $\Prb(Z>u)=\Prb(Z<v)$, and hence
\[
 \Prb(X<Y)-\Prb(X>Y)=\Prb(u\le Z<v)=Q(u,v).
\]
It follows from~\eqref{eq:prob} that $W_2(u,v)=D(u,v)$ and
\begin{equation}\label{eq:H2}
 H_2(n,n)=\sum_{0\le u<v\le N}D(u,v)^2.
\end{equation}

For any three real numbers, including ties, the identity
\begin{equation}\label{eq:sign3}
 \sgn((y-x)(z-x)(z-y))
 =\sgn(y-x)-\sgn(z-x)+\sgn(z-y)
\end{equation}
follows by checking their relative order. Taking expectations gives
\begin{equation}\label{eq:w3}
 w_3(u,v,w)=Q(u,v)-Q(u,w)+Q(v,w)
 \quad(0\le u<v<w).
\end{equation}
Thus, if
\[
 A(u,v,w)=2^wD(u,v)-2^vD(u,w)+2^uD(v,w),
\]
then $W_3(u,v,w)=A(u,v,w)$ and
\begin{equation}\label{eq:H3}
 H_3(n,n)=\sum_{0\le u<v<w\le N}A(u,v,w)^2.
\end{equation}
Equations~\eqref{eq:H2} and~\eqref{eq:H3} are exact integer identities.

\section{A uniform local expansion}

Write $u=N-a$, $v=N-b$, where $0\le b<a\le L$ and
$L=O(\log N)$, and put $d=a-b$, $s=a+b$.

\begin{lemma}\label{lem:local}
Uniformly for these offsets,
\begin{equation}\label{eq:local}
 Q(N-a,N-b)=\frac{d}{\sqrt{\pi N}}
 \left[1+\frac{s/4-d^2/12-7/24}{N}\right]
 +O\left(\frac{(1+L)^5}{N^{5/2}}\right).
\end{equation}
\end{lemma}

\begin{proof}
Put $m=2N-s$. The summation indices in~\eqref{eq:D} can be written
$j=m/2+t$, with $t=-d/2,-d/2+1,\ldots,d/2-1$.
The numbers $t$ may be half-integers, but $j$ and $m-j$ are integers.
The Stirling expansion with remainder gives, uniformly here,
\begin{equation}\label{eq:mass}
 2^{-m}\binom m{m/2+t}
 =\sqrt{\frac2{\pi m}}
 \left[1-\frac{2t^2+1/4}{m}
 +O\left(\frac{(1+|t|)^4}{m^2}\right)\right].
\end{equation}
For completeness, set $x=2t/m$. Applying the logarithmic Stirling
formula to $m!$, $(m/2+t)!$, and $(m/2-t)!$ gives
\begin{align*}
 \log\!\left(2^{-m}\binom m{m/2+t}\sqrt{\frac{\pi m}{2}}\right)
 ={}&-\frac m2\bigl[(1+x)\log(1+x)+(1-x)\log(1-x)\bigr]\\
 &-\frac12\log(1-x^2)+\frac1{12m}
 -\frac1{12(m/2+t)}-\frac1{12(m/2-t)}+O(m^{-3}).
\end{align*}
Since $|t|=O(\log N)$ and $m\asymp N$, the three factorial arguments
are uniformly comparable to $N$, and eventually $|x|\le1/2$.
Taylor expansion therefore makes the last display equal to
\[
 -\frac{2t^2+1/4}{m}
 +O\left(\frac{t^2}{m^2}+\frac{t^4+1}{m^3}\right).
\]
Exponentiating proves~\eqref{eq:mass}, with an error constant uniform
over the logarithmic window.
Substituting $m=2N-s$ into~\eqref{eq:mass} gives each mass as
\[
 \frac1{\sqrt{\pi N}}
 \left[1+\frac{s/4-t^2-1/8}{N}
 +O\left(\frac{(1+L)^4}{N^2}\right)\right].
\]
Finally,
\[
 \sum_{k=0}^{d-1}(k-d/2)^2=\frac{d(d^2+2)}{12}.
\]
There are $d\le L$ terms, so summing yields~\eqref{eq:local}.
\end{proof}

For $r=2$, the lemma immediately implies
\begin{equation}\label{eq:w2local}
 w_2(N-a,N-b)=\frac{a-b}{\sqrt{\pi N}}(1+o(1))
\end{equation}
uniformly in the window. For $r=3$, take $a>b>c\ge0$ and apply the
lemma to the three terms in~\eqref{eq:w3}. The coefficients of
$N^{-1/2}$ cancel because
$(a-b)-(a-c)+(b-c)=0$. In the next coefficient, the terms involving
$ds=a^2-b^2$ and the constant multiple of $d$ also cancel. Since
\[
 (a-b)^3-(a-c)^3+(b-c)^3
 =-3(a-b)(a-c)(b-c),
\]
we obtain, with $\Delta(a,b,c)=(a-b)(a-c)(b-c)$,
\begin{equation}\label{eq:w3local}
 w_3(N-a,N-b,N-c)=
 \frac{\Delta(a,b,c)}{4\sqrt\pi\,N^{3/2}}
 +O\left(\frac{(1+L)^5}{N^{5/2}}\right).
\end{equation}
For distinct integer offsets $\Delta(a,b,c)\ge2$, so the relative
error in~\eqref{eq:w3local} tends uniformly to zero. The same is true
in~\eqref{eq:w2local}, since $a-b\ge1$. The next-order term in the
local binomial expansion is essential in the three-row case.

\section{Tail control and evaluation of the constants}

Let $r\in\{2,3\}$ and $R=\binom r2$. Replacing the row indices in
\eqref{eq:squares} by decreasing offsets gives
\begin{equation}\label{eq:normalized}
 \frac{N^R H_r(n,n)}{2^{2rN}}
 =\sum_{N\ge a_1>\cdots>a_r\ge0}
 4^{-\sum_i a_i}N^R
 w_r(N-a_1,\ldots,N-a_r)^2.
\end{equation}
Choose $L=\lceil6\log_4N\rceil$. By~\eqref{eq:globalbound}, the
part of the sum with $a_1>L$ is at most
\begin{align}
 N^R\sum_{\max_i a_i>L}4^{-\sum_i a_i}
 &\le N^R r\left(\sum_{k=L+1}^\infty4^{-k}\right)
                \left(\sum_{k=0}^\infty4^{-k}\right)^{r-1}\notag\\
 &=\frac r3\left(\frac43\right)^{r-1}N^R4^{-L}
 =O(N^{R-6})=o(1).\label{eq:tail}
\end{align}
For this upper bound we have enlarged the sum to all nonnegative
integer tuples.

Inside the window,~\eqref{eq:w2local} and~\eqref{eq:w3local} have
uniform relative errors. A squared Vandermonde is a fixed polynomial
in the offsets, so its sum against $4^{-\sum_i a_i}$ converges.
Consequently the error in the weighted window sum tends to zero,
while the truncated main sum tends to the full convergent sum.
Combining this fact with~\eqref{eq:tail} proves
\begin{equation}\label{eq:limits}
 H_2(n,n)\sim\frac{16^N}{\pi N}Z_2(1/4),\qquad
 H_3(n,n)\sim\frac{64^N}{16\pi N^3}Z_3(1/4),
\end{equation}
where, for $0<q<1$,
\[
 Z_r(q)=\sum_{0\le b_1<\cdots<b_r}
       \prod_{i<j}(b_j-b_i)^2q^{b_1+\cdots+b_r}.
\]
This argument explicitly bounds the whole tail, rather than
interchanging a pointwise limit and an unbounded sum.

Put $\mu_k(q)=\sum_{b\ge0}b^kq^b$, with $b^0=1$ also for $b=0$.
Equivalently,
\[
 \mu_k(q)=\left(q\frac{d}{dq}\right)^k\frac1{1-q}.
\]
Cauchy--Binet for the weighted Vandermonde matrix, first with bounded
column indices and then in the limit, gives
\[
 Z_r(q)=\det[\mu_{i+j}(q)]_{0\le i,j<r}.
\]
Every moment converges for $0<q<1$, and a finite determinant is a
polynomial in its entries, which justifies that limit. Direct
differentiation and the two small determinants give
\[
 Z_2(q)=\frac q{(1-q)^4},\qquad
 Z_3(q)=\frac{4q^3}{(1-q)^9}.
\]
In particular,
\[
 Z_2(1/4)=\frac{64}{81},\qquad
 Z_3(1/4)=\frac{16384}{19683}.
\]
Finally $N=n-2$ and $N/n\to1$, while
\[
 \frac{64/81}{16^2}=\frac1{324},\qquad
 \frac{16384/19683}{16\cdot64^2}=\frac1{78732}.
\]
Substitution into~\eqref{eq:limits} proves Theorem~\ref{thm:main}.
\qed

\section{Reproducible finite checks}

The accompanying script \texttt{hardinian.py} evaluates the exact
integer formulas through $n=201$. It also enumerates arrays directly
from the definition for $1\le n\le6$ and evaluates the complementary
minors independently for $3\le n\le8$, in both cases for $r=2,3$.
All comparisons pass. The computed values agree with the 37 supplied
terms of OEIS A253217 and the first 24 supplied terms of A252998,
together with its entries at $n=100,101$~\cite{oeis2,oeis3}.
The script includes exact symbolic checks of $Z_2$ and $Z_3$ when
SymPy is installed. The accompanying result file records an actual
successful run. These checks help detect indexing and constant errors;
the limits for unbounded $n$ are proved in Sections~4--5.

\paragraph{AI assistance and verification status.}
This research was completed with AI assistance. The mathematical
argument, source references, and reproducible checks are provided for
evaluation. This version does not claim proof-assistant verification
or external peer review.

\begin{thebibliography}{9}
\bibitem{dk}
R.~Dougherty-Bliss and M.~Kauers,
\emph{Hardinian Arrays},
Electronic Journal of Combinatorics \textbf{31}(2) (2024), P2.9.
\href{https://doi.org/10.37236/12358}{doi:10.37236/12358}.

\bibitem{oeis2}
The OEIS Foundation Inc.,
\emph{The On-Line Encyclopedia of Integer Sequences},
\href{https://oeis.org/A253217}{A253217};
\href{https://oeis.org/A253217/b253217.txt}{table of values}.
Accessed October 13, 2026.

\bibitem{oeis3}
The OEIS Foundation Inc.,
\emph{The On-Line Encyclopedia of Integer Sequences},
\href{https://oeis.org/A252998}{A252998};
\href{https://oeis.org/A252998/b252998.txt}{table of values}.
Accessed October 13, 2026.
\end{thebibliography}
\end{document}
